\section{Results}
\subsection{Auditing one attribute misses the intersection; auditing both marginals does not}
\label{sec:mask}
On natural BRFSS at $\lambda{=}8$ (Table~\ref{tab:mask}; 5 natural and 20 capped seeds) the sex-only model cuts the sex DP gap by 87\% (0.112$\to$0.015) but the age$\times$sex gaps by only 33\% (DP 0.254$\to$0.171) and 28\% (EO 0.459$\to$0.333), at AUROC 0.798$\to$0.789. A sex-only audit would pass this model on DP (0.015) although its intersectional DP gap is 11$\times$ larger; capped data agree (DP$_\text{int}$ 0.251$\to$0.167, EO$_\text{int}$ 0.453$\to$0.325). The masking is visible on the age marginal itself: the sex-only model's age DP gap is 0.154 against 0.152 for FedAvg (natural; 0.151 against 0.150 capped), whereas the marginal-both model reaches 0.024 (0.034 capped), so the decision rule fixed in advance (Sec.~III-A) is met in both scenarios. Penalizing the sex and age marginals together closes the intersectional gap almost as well as the four cells: DP$_\text{int}$ 0.031 vs.\ 0.032 and EO$_\text{int}$ 0.232 vs.\ 0.228 (natural), 0.043 vs.\ 0.048 and 0.182 vs.\ 0.171 (capped). With two binary attributes the masking is therefore a failure of auditing \emph{one} attribute, not of marginal fairness; whether crossing matters with more or finer attributes is untested.

\begin{table}[t]
\centering
\caption{Masking at $\lambda{=}8$ (fine-grid suite; logistic regression, rule (q)). Reductions are relative to FedAvg. ``Marginal'' penalizes the sex and age indicators together; ``Intersectional'' the four age$\times$sex cells. DP$_\text{age}$: old vs.\ young selection-rate gap.}
\label{tab:mask}
\scriptsize
\setlength{\tabcolsep}{3pt}
\begin{tabular}{lccccccc}
\toprule
Penalty & AUROC & DP$_\text{sex}$ & DP$_\text{age}$ & DP$_\text{int}$ & $\Delta$DP$_\text{int}$ & EO$_\text{int}$ & $\Delta$EO$_\text{int}$\\
\midrule
\input{tab3_mask.tex}
\bottomrule
\end{tabular}
\end{table}

\subsection{What penalties cost, how strength matters, and group thresholds}
\label{sec:post}
\textbf{Cost.} Intersectional penalties close most of the DP gap, but AUROC falls from 0.798 to 0.680 (natural, $\lambda{=}8$, 20 seeds; 0.797 to 0.681 capped) and sensitivity at the deployable rule (fx) from 0.310 to 0.210 (0.215 capped). The cost is not due to federation: the penalized centralized control reaches 0.679 (natural) and 0.678 (capped), and unpenalized FedAvg and its centralized control are identical (0.798).

\textbf{Strength.} The EO gap is U-shaped in $\lambda$. For the intersectional penalty (natural) EO$_\text{int}$ is 0.459 at $\lambda{=}0$, 0.371 at 0.1, 0.100 at 0.5, 0.105 at 1, 0.218 at 4 and 0.228 at 8 (capped: 0.453, 0.352, 0.129, 0.126, 0.165, 0.171), and at the EO optimum ($\lambda{=}0.5$) the cost is smaller: AUROC 0.751 against 0.797 for FedAvg. On 5 natural seeds the age coefficient is $+0.766$ at $\lambda{=}0$, $-0.001$ at 0.5 and $-0.24$ at 8, so the penalty first removes the model's dependence on age and then over-corrects (a correlational explanation, not a causal test). Comparing mitigations at one large $\lambda$ therefore overstates the penalty's EO gap.

\textbf{Group thresholds} (10 natural / 20 capped seeds, rule (fx); the frontier offsets elsewhere use 20 seeds and rule (q), so numbers differ slightly, e.g.\ naive $\lambda{=}1$ AUROC 0.716 vs.\ 0.717). Equal-TPR thresholds fitted on FedAvg training scores keep AUROC (0.797) and give EO$_\text{int}$ 0.087 (natural) and 0.142 (capped) under (fx) (Table~\ref{tab:thr}). Against the penalty, natural paired differences are $-0.015$ [$-0.031,+0.000$] at $\lambda{=}1$, $-0.133$ [$-0.144,-0.122$] at 4 and $-0.145$ [$-0.155,-0.135$] at 8; capped, $+0.018$ [$-0.016,+0.049$] at 1 and $-0.026$ [$-0.061,+0.009$] at 8. Thresholds thus beat a strong penalty on EO in natural data only and tie a weak one, but the weak penalty already loses sensitivity (0.231 vs.\ 0.310) and PPV (0.233 vs.\ 0.270). Fitted on a validation split of about 30 respondents per client (capped), EO is worse (0.237) because few positives determine four thresholds; natural data match (0.090 vs.\ 0.087). Against a label-conditional EO penalty, natural thresholds are lower at every strength tested ($-0.009$ [$-0.012,-0.006$] at $\lambda{=}0.5$; $-0.107$ [$-0.123,-0.090$] at 4) with 0.067--0.112 more AUROC, and capped data are indistinguishable; high-$\lambda$ EO-penalty numbers are unreliable (two unguarded fits collapsed, \sref{app:conv}). Thresholds need the sensitive attribute at decision time, and fitting them in FL would require group-wise score statistics from very small cells.

\begin{table}[t]
\centering
\caption{Deployable rule: equal-TPR thresholds on FedAvg versus the intersectional penalty at several $\lambda$. Mean over 10 natural / 20 capped seeds.}
\label{tab:thr}
\scriptsize
\setlength{\tabcolsep}{3pt}
\begin{tabular}{lccccc}
\toprule
Method & AUROC & DP$_\text{int}$ & EO$_\text{int}$ & Sens. & PPV\\
\midrule
\input{tab3_thr.tex}
\bottomrule
\end{tabular}
\end{table}

Group-wise, the intersectional penalty lowers AUROC in every group (older women/men 0.72/0.70 to 0.67/0.65; younger 0.80/0.79 to 0.78/0.76) and shifts calibration-in-the-large from about zero to $+0.065$ (younger women), $+0.061$ (younger men), $-0.004$ (older women) and $-0.084$ (older men), i.e.\ it under-predicts risk for older men by 8 points (\sref{tab:grp}). All gaps are measured at one operating point. FedAvg's EO$_\text{int}$ grows from 0.307 (5\% selected) to 0.698 (20\%) in natural data and the youngest group of women has TPR 0.000 at 5, 9.3 and 15\%; comparisons between mitigations change with the rate (\sref{tab:op}), so signs and sizes of EO effects are specific to the operating points used.

\subsection{Oversampling: what it does, and what drives it}
\label{sec:over}
On the dense grid, local oversampling lowers DP ($-0.012$ [$-0.014,-0.010$]) and raises EO ($+0.059$ [$+0.054,+0.064$]) at equal AUROC in natural data (per-seed range $+0.044$ to $+0.081$; Holm-adjusted $p<0.001$; Fig.~\ref{fig:dense}). In capped data neither is distinguishable from zero: DP $-0.006$ [$-0.014,+0.002$] (adjusted $p=0.47$), EO $+0.020$ [$-0.004,+0.041$] (adjusted $p=0.45$); among per-$\lambda$ equal-AUROC offsets only $\lambda{=}8$ excludes zero, and at $\lambda{=}0.5$ the offset is not estimable because the oversampled models lie above the naive frontier in 17 of 20 seeds (\sref{tab:perlam}); the equal-$\lambda$ comparison does show a significant capped EO difference at $\lambda{=}0.5$ ($+0.033$ [$+0.015,+0.052$], Table~\ref{tab:eqlam}), the only capped point that does, which the average over $\lambda$ dilutes. The comparator matters: at equal $\lambda$ the oversampled natural models have \emph{lower} EO than the naive penalty for $\lambda\ge4$ because they keep 0.02--0.03 more AUROC, and higher EO at $\lambda\le1$ (Table~\ref{tab:eqlam}), so ``raises EO'' holds at equal AUROC, not at equal penalty strength.

\begin{figure*}[t]
\centering
\includegraphics[width=0.8\textwidth]{../figs/frontier_dense_natural.pdf}
\caption{Natural data, dense grid $\lambda\in\{0.5,1,2,4,8,16\}$, mean of 20 seeds. The naive EO curve is U-shaped (minimum near $\lambda{=}0.5$--1), so a chord between $\lambda{=}1$ and 4 overstates the oversampling offset.}
\label{fig:dense}
\end{figure*}

\begin{table}[t]
\centering
\caption{Equal-$\lambda$ comparison, oversample vs.\ naive intersectional penalty, logistic regression, 20 seeds. $\Delta$EO is the paired difference [95\% CI].}
\label{tab:eqlam}
\footnotesize
\resizebox{\columnwidth}{!}{\input{tab5_eqlam.tex}}
\end{table}

Balancing \emph{labels} by duplication alone reproduces about half (49\%) of the natural dense EO shift ($+0.029$ [$+0.026,+0.032$]); the paired difference (oversample $-$ label-balance) is $+0.030$ [$+0.027,+0.034$], so the other half is not explained by label balance. In capped data both shifts are indistinguishable from zero ($+0.020$ and $+0.018$; paired difference $+0.001$ [$-0.011,+0.015$]), so no share can be given. Re-weighting by \emph{group alone} lowers EO by 0.011 (natural and capped; adjusted $p<0.001$ and 0.003) with DP $+0.001$ to $+0.002$. Oversampling raises ECE from 0.006 to 0.351 (natural) because the training prevalence moves toward 50\%; an intercept shift removes this without changing DP or EO, which depend only on score ranks.

\subsection{Is it specific to federation? Client skew}
\label{sec:fedspec}
\begin{table*}[t]
\centering
\caption{Centralized matched control vs.\ federated: frontier offsets, logistic regression, rule~(q), mean [95\% CI over seeds]. ``Diff'' = federated $-$ central (paired); 20 seeds per scenario; natural: dense frontier $\lambda\in\{0.5,1,2,4,8,16\}$, capped: core grid.}
\label{tab:cvf}
\small
\resizebox{\textwidth}{!}{%
\begin{tabular}{lcccccc}
\toprule
& \multicolumn{3}{c}{Natural} & \multicolumn{3}{c}{Capped}\\
\cmidrule(lr){2-4}\cmidrule(lr){5-7}
Variant & Central & Federated & Diff & Central & Federated & Diff\\
\midrule
\input{tab2_cvf.tex}
\bottomrule
\end{tabular}}
\end{table*}
\begin{table*}[t]
\centering
\caption{Dirichlet-skewed client partitions of the same training data (10 seeds): offsets vs.\ the naive penalty on the core-grid chord (which overstates natural EO offsets, Sec.~III-C and Fig.~\ref{fig:dense}), rule~(q), mean [95\% CI]. Federated, centralized, and their paired difference. Exploratory.}
\label{tab:dir}
\scriptsize
\resizebox{\textwidth}{!}{%
\begin{tabular}{llcccccc}
\toprule
& & \multicolumn{3}{c}{DP$_\text{int}$} & \multicolumn{3}{c}{EO$_\text{int}$}\\
\cmidrule(lr){3-5}\cmidrule(lr){6-8}
Partition & Variant & Fed. & Central & Diff & Fed. & Central & Diff\\
\midrule
\input{tab2_dir.tex}
\bottomrule
\end{tabular}}
\end{table*}

A centralized control with the same number of gradient steps (Table~\ref{tab:cvf}) gives oversampling DP $-0.016$ and EO $+0.052$ (natural); federation adds EO $+0.008$ [$+0.005,+0.010$] and nothing measurable in capped data (EO $+0.009$ [$-0.014,+0.031$]), so on near-IID clients the effect is a property of the mitigation, not of federation. Under synthetic Dirichlet client skew ($\alpha{=}0.3$; Table~\ref{tab:dir}; offsets there use the core-grid chord, so only the federated $-$ central differences are interpreted) that changes: oversampling's EO effect is smaller federated than centrally (federated $-$ central $-0.055$ [$-0.077,-0.032$]), and group re-weighting with local counts turns harmful (difference $+0.032$ [$+0.014,+0.053$]), because a client missing a group assigns it a very different weight than the pooled data would; weights from four pooled counts repair it (EO offset $-0.008$ [$-0.012,-0.005$], \sref{tab:glob}). Further checks (leave-states-out, client size, age cut-off, shrinkage, an MLP, an exponentiated-gradient reference, UCI, convergence) are in the supplement (\srefr{app:robust}{sec:mlp}).
